\section{Функціональні вимоги}
\label{sec:funktsionalni-vymohy}

\subsection{Автентифікація та управління акаунтом}

\begin{table}[h]
\caption{Вимоги до автентифікації та акаунту}
\label{tab:auth}
\begin{tabular}{L{2.5cm} L{11.5cm}}
\toprule
\textbf{ID} & \textbf{Вимога} \\
\midrule
REQ-F-001 & Система повинна надавати можливість реєстрації за email, паролем та нікнеймом (усі три поля~--- обов'язкові). Повторна реєстрація з тим самим email неможлива. \sked{K7} \\
\addlinespace
REQ-F-002 & Система повинна надавати вхід до акаунту за email та паролем. \sked{K8} \\
\addlinespace
REQ-F-003 & Система повинна надавати можливість виходу з акаунту. Після виходу доступ до захищених сторінок вимагає повторної авторизації. \sked{K9} \\
\addlinespace
REQ-F-029 & Користувач може змінити нікнейм у налаштуваннях акаунту. \sked{K49} \\
\bottomrule
\end{tabular}
\end{table}

\subsection{Каталог}

\begin{table}[h]
\caption{Вимоги до каталогу}
\label{tab:kataloh}
\begin{tabular}{L{2.5cm} L{11.5cm}}
\toprule
\textbf{ID} & \textbf{Вимога} \\
\midrule
REQ-F-004 & Система повинна відображати особистий каталог~--- список усіх елементів контенту користувача. \sked{K10} \\
\addlinespace
REQ-F-005 & Система повинна надавати фільтрування каталогу за типом елемента (книга~/ серіал~/ фільм). \sked{K11} \\
\addlinespace
REQ-F-006 & Система повинна надавати фільтрування каталогу за статусом елемента. Доступні значення статусу: «заплановано», «у процесі», «завершено», «відкладено». \sked{K12} \\
\addlinespace
REQ-F-023 & Система повинна надавати пошук елементів каталогу за назвою. \sked{K29} \\
\addlinespace
REQ-F-026 & Система повинна надавати сортування каталогу: за датою додавання (новіші~/ старіші) та за оцінкою (від вищої до нижчої і навпаки). \sked{K32} \\
\bottomrule
\end{tabular}
\end{table}

\subsection{Управління елементами контенту}

\begin{table}[h]
\caption{Вимоги до управління елементами контенту}
\label{tab:content}
\begin{tabular}{L{2.5cm} L{11.5cm}}
\toprule
\textbf{ID} & \textbf{Вимога} \\
\midrule
REQ-F-007 & Система повинна надавати можливість додавання нового елемента контенту з атрибутами: назва (обов'язково), тип (обов'язково), статус, оцінка~(0–10), нотатки, теги. Форма відображає тип-специфічні необов'язкові поля (REQ-F-024) та поле обкладинки (REQ-F-025). \sked{K13} \\
\addlinespace
REQ-F-008 & Система повинна надавати можливість редагування будь-якого атрибута наявного елемента контенту. \sked{K14} \\
\addlinespace
REQ-F-009 & Система повинна надавати можливість видалення елемента контенту з каталогу. Видалений елемент автоматично зникає з усіх добірок. \sked{K15} \\
\addlinespace
REQ-F-024 & Форма додавання та редагування елемента повинна відображати тип-специфічні необов'язкові поля залежно від обраного типу (см.~\tabref{tab:type-fields}). \sked{K30} \\
\addlinespace
REQ-F-025 & Кожен елемент контенту повинен мати обкладинку: за замовчуванням~--- одна з трьох ілюстрацій відповідно до типу. Користувач може замінити її власним зображенням (необов'язково). \sked{K31} \\
\addlinespace
REQ-F-028 & Формати зображення для обкладинки~--- JPEG, PNG, WebP; максимальний розмір файлу~--- 5~МБ. Система повідомляє користувача при порушенні обмежень. \sked{K48} \\
\bottomrule
\end{tabular}
\end{table}

\begin{table}[h]
\caption{Тип-специфічні поля елементів контенту}
\label{tab:type-fields}
\begin{tabular}{L{3cm} L{11cm}}
\toprule
\textbf{Тип} & \textbf{Необов'язкові поля} \\
\midrule
Книга & Автор, кількість сторінок \\
Серіал & Кількість сезонів \\
Фільм & Режисер, тривалість \\
\bottomrule
\end{tabular}
\end{table}
\sked{K30}

\subsection{Теги}

\begin{table}[h]
\caption{Вимоги до тегів}
\label{tab:tehy}
\begin{tabular}{L{2.5cm} L{11.5cm}}
\toprule
\textbf{ID} & \textbf{Вимога} \\
\midrule
REQ-F-010 & Система повинна надавати можливість додавання довільних тегів до елемента контенту. При введенні тегу система пропонує автодоповнення на основі раніше введених тегів користувача. Один елемент може мати кілька тегів. Тег можна видалити в будь-який момент. \sked{K16} \\
\addlinespace
REQ-F-011 & Система повинна надавати фільтрування та пошук по каталогу за тегом. \sked{K17} \\
\bottomrule
\end{tabular}
\end{table}

\subsection{Sessions (повторні перегляди~/ прочитання)}

\begin{table}[h]
\caption{Вимоги до Sessions}
\label{tab:sessions}
\begin{tabular}{L{2.5cm} L{11.5cm}}
\toprule
\textbf{ID} & \textbf{Вимога} \\
\midrule
REQ-F-012 & Система повинна надавати можливість фіксації кількох Sessions для одного елемента контенту; кожна Session має власні дату, оцінку та нотатку. \sked{K18} \\
\addlinespace
REQ-F-013 & Система повинна відображати список усіх Sessions елемента контенту у зворотному хронологічному порядку (від найновішої до найстарішої), із зазначенням загальної кількості. \sked{K19} \\
\addlinespace
REQ-F-014 & Система повинна надавати можливість редагування та видалення окремої Session. \sked{K20} \\
\bottomrule
\end{tabular}
\end{table}

\subsection{Добірки}

\begin{table}[h]
\caption{Вимоги до добірок}
\label{tab:dobirkы}
\begin{tabular}{L{2.5cm} L{11.5cm}}
\toprule
\textbf{ID} & \textbf{Вимога} \\
\midrule
REQ-F-015 & Система повинна надавати можливість додавання одного елемента каталогу до кількох добірок одночасно (зв'язок many-to-many). \sked{K21} \\
\addlinespace
REQ-F-016 & Система повинна надавати можливість створення нової іменованої добірки. \sked{K22} \\
\addlinespace
REQ-F-017 & Система повинна надавати можливість редагування назви добірки. \sked{K23} \\
\addlinespace
REQ-F-018 & Система повинна надавати можливість видалення добірки. Видалення добірки не видаляє елементи з каталогу. \sked{K24} \\
\addlinespace
REQ-F-019 & Система повинна надавати можливість видалення елемента з добірки без видалення його з каталогу чи з інших добірок. \sked{K25} \\
\addlinespace
REQ-F-020 & Система повинна відображати вміст добірки~--- список елементів, що до неї входять, з назвою, типом і статусом кожного. \sked{K26} \\
\addlinespace
REQ-F-027 & Система повинна надавати сортування елементів всередині добірки: за датою додавання (новіші~/ старіші) та за оцінкою (від вищої до нижчої і навпаки). \sked{K33} \\
\bottomrule
\end{tabular}
\end{table}

\subsection{Автозаповнення дати завершення}

\begin{table}[h]
\caption{Вимога до дати завершення}
\label{tab:data}
\begin{tabular}{L{2.5cm} L{11.5cm}}
\toprule
\textbf{ID} & \textbf{Вимога} \\
\midrule
REQ-F-021 & При зміні статусу елемента на «завершено» поле дати завершення повинно автоматично заповнюватися поточною датою, якщо дату не введено вручну. \sked{K27} \\
\bottomrule
\end{tabular}
\end{table}

\subsection{Особиста статистика}

\begin{table}[h]
\caption{Вимоги до особистої статистики}
\label{tab:stat}
\begin{tabular}{L{2.5cm} L{11.5cm}}
\toprule
\textbf{ID} & \textbf{Вимога} \\
\midrule
REQ-F-022 & Система повинна відображати особисту статистику: кількість елементів зі статусом «завершено» у розрізі типів (книги~/ серіали~/ фільми) та за обраний часовий проміжок. \sked{K28} \\
\addlinespace
REQ-F-022a & Вибір часового проміжку для статистики~--- з наперед визначених варіантів: «цей місяць», «цей рік», «весь час». \sked{K51} \\
\bottomrule
\end{tabular}
\end{table}
